% =====================================================================
%  Systematic Review manuscript template (PRISMA 2020) -- engineering
%  Compiles with: pdflatex main && bibtex main && pdflatex main && pdflatex main
%
%  SWITCHING TO A VENUE CLASS
%   IEEE journal:    \documentclass[journal]{IEEEtran}      bib style: IEEEtran
%   IEEE conference: \documentclass[conference]{IEEEtran}   bib style: IEEEtran
%   Elsevier:        \documentclass[preprint,12pt]{elsarticle} (+ frontmatter env)
%                    bib style: elsarticle-num  or elsarticle-harv
%   Springer Nature: \documentclass[sn-mathphys-num]{sn-jnl}
%   ASME:            \documentclass{asmejour}
%   ACM:             \documentclass[acmsmall]{acmart}       bib style: ACM-Reference-Format
%   MDPI:            use the journal's Definitions/mdpi.cls template
%  Then replace the title/author block below with the class's own macros.
% =====================================================================
\documentclass[11pt,a4paper]{article}
\usepackage[T1]{fontenc}
\usepackage[utf8]{inputenc}
\usepackage{lmodern}
\usepackage[margin=2.5cm]{geometry}
\usepackage{amsmath,amssymb}
\usepackage{graphicx}
\usepackage{booktabs,tabularx,longtable,multirow}
\usepackage{siunitx}
\usepackage[table]{xcolor}
\usepackage{tikz}
\usetikzlibrary{arrows.meta,positioning}
\usepackage{caption,subcaption}
\usepackage[numbers,sort&compress]{natbib}
\usepackage[hidelinks]{hyperref}
\usepackage{cleveref}

\newcommand{\authorinput}[1]{\textcolor{red}{[AUTHOR INPUT: #1]}}
\newcommand{\verify}[1]{\textcolor{orange}{[VERIFY: #1]}}

\title{<Topic>: A Systematic Review of <Scope> and Research Agenda}
\author{\authorinput{Authors, affiliations, ORCID, corresponding e-mail}}
\date{}

\begin{document}
\maketitle

\begin{abstract}
% 150--250 words. Context -> gap -> objective -> methods (PRISMA 2020; databases; date of last
% search; eligibility; N screened; N included; appraisal) -> quantitative key findings -> implications
% -> agenda. No citations. Fill every number from prisma_counts.json.
\end{abstract}

\noindent\textbf{Keywords:} systematic literature review; PRISMA; <kw3>; <kw4>; <kw5>

% Elsevier: Highlights (3-5 bullets, <= 85 characters each) go in a separate file.

%======================================================================
\section{Introduction}\label{sec:intro}
% 1. Engineering context and stakes (quantified, cited)
% 2. Brief technical background
% 3. Existing reviews and their limitations -> Table~\ref{tab:prior}
% 4. Objectives and RQs
% 5. Contributions (3--5 bullets)
% 6. Paper organisation

\begin{table}[t]
\centering
\caption{Comparison with existing reviews on <topic>.}\label{tab:prior}
\small
\begin{tabularx}{\linewidth}{@{}l c X c c X@{}}
\toprule
Review & Year & Scope & Systematic? & \# studies & Limitation addressed here \\
\midrule
\cite{key} & 20XX & & & & \\
This review & 20XX & & PRISMA 2020 & & \\
\bottomrule
\end{tabularx}
\end{table}

The main contributions of this review are:
\begin{itemize}
  \item a PRISMA 2020-compliant synthesis of $N$ studies published between YYYY and YYYY;
  \item a taxonomy of \ldots (\cref{fig:taxonomy});
  \item a critical appraisal of validation maturity and reproducibility;
  \item an evidence-based research agenda (\cref{tab:agenda}).
\end{itemize}

%======================================================================
\section{Background}\label{sec:background}
% Definitions and fundamentals only (journal versions).

%======================================================================
\section{Methodology}\label{sec:method}
This review follows the PRISMA 2020 statement~\cite{page2021prisma}% verify key in references.bib
\ and, where applicable, the guidelines of Kitchenham and Charters for systematic reviews in engineering disciplines. The protocol was registered on \authorinput{OSF Registries, DOI} before screening began.

\subsection{Research questions}
\begin{table}[h]
\centering\small
\caption{PICOC framework and research questions.}\label{tab:picoc}
\begin{tabularx}{\linewidth}{@{}l X@{}}
\toprule
Population & \\ Intervention & \\ Comparison & \\ Outcome & \\ Context & \\
\midrule
RQ1 & \\ RQ2 & \\ RQ3 & \\ RQ4 & \\
\bottomrule
\end{tabularx}
\end{table}

\subsection{Eligibility criteria}
% Table of IC/EC with IDs.

\subsection{Information sources and search strategy}
% Databases + interface + date of last search + hits (table). Main string here; all strings in Supplementary S2.

\subsection{Selection process}
% Two stages; reviewers; independence; AI assistance disclosure; kappa on the verification sample.

\subsection{Data collection and data items}
% Extraction form fields (Supplementary S4), piloting, verification.

\subsection{Quality assessment}
% Instrument, items, scoring, thresholds, use in synthesis.

\subsection{Synthesis methods}
% SWiM narrative / taxonomy / meta-analysis (model, heterogeneity, software); sensitivity analyses; certainty (adapted GRADE).

%======================================================================
\section{Results}\label{sec:results}
\subsection{Study selection}
\begin{figure}[t]
\centering
\resizebox{\linewidth}{!}{\input{prisma_flow.tex}}
\caption{PRISMA 2020 flow diagram of the study selection process.}\label{fig:prisma}
\end{figure}

\subsection{Characteristics of included studies}
% Publications per year, venues, evidence types; summary table (condensed; full in Supplementary).

\subsection{Quality of included studies}

\subsection{RQ1: <title>}
% Organise by taxonomy, not paper-by-paper. Claim -> k studies [refs] -> magnitude -> conditions -> explanation -> certainty.

\subsection{RQ2: <title>}
\subsection{RQ3: <title>}

\begin{table}[t]
\centering\small
\caption{Summary of findings and certainty of evidence.}\label{tab:sof}
\begin{tabularx}{\linewidth}{@{}X c l l X@{}}
\toprule
Finding & $k$ & Evidence type & Certainty & Rationale \\
\midrule
 & & & & \\
\bottomrule
\end{tabularx}
\end{table}

%======================================================================
\section{Discussion}\label{sec:discussion}
\subsection{Principal findings}
\subsection{Comparison with prior reviews}
\subsection{Proposed framework}
% \begin{figure}... \label{fig:taxonomy}
\subsection{Implications for engineering practice}
\subsection{Research agenda}
\begin{table}[t]
\centering\small
\caption{Research agenda derived from the identified gaps.}\label{tab:agenda}
\begin{tabularx}{\linewidth}{@{}X X X X@{}}
\toprule
Gap & Evidence & Why it matters & Suggested direction \\
\midrule
 & & & \\
\bottomrule
\end{tabularx}
\end{table}
\subsection{Limitations and threats to validity}
% Of the evidence (23b) and of the review process (23c): databases, language, period, AI-assisted screening, publication bias.

%======================================================================
\section{Conclusions}\label{sec:conclusion}

%======================================================================
\section*{CRediT authorship contribution statement}
\authorinput{CRediT roles}

\section*{Declaration of generative AI and AI-assisted technologies}
During the preparation of this work the author(s) used \authorinput{tool, version} in order to \authorinput{assist title/abstract screening and data extraction; language editing}. All AI-assisted decisions were verified as described in \cref{sec:method}. After using this tool, the author(s) reviewed and edited the content as needed and take(s) full responsibility for the content of the publication.

\section*{Declaration of competing interest}
\authorinput{statement}

\section*{Funding}
\authorinput{funding sources, grant numbers}

\section*{Data availability}
Search strings, screening logs, the data-extraction sheet, quality scores and analysis scripts are available at \authorinput{OSF/Zenodo DOI}.

\bibliographystyle{unsrtnat}
\bibliography{references}

\appendix
\section*{Supplementary material}
S1 PRISMA 2020 checklist; S2 full search strings; S3 screening logs; S4 data-extraction table; S5 quality-assessment scores.

\end{document}
